% ============================================================
% conference-poster-beamer.tex — math-heavy conference poster
% (beamerposter route; template credits: the beamerposter package's
% own sample, restyled to the pack's five-color system)
%
% Copy in full, then edit content. Placeholder text doubles as
% filling instructions.
%
% Compile:  xelatex conference-poster-beamer.tex   (x2 for refs)
% Requires: beamerposter.sty — on TeX Live:
%             tlmgr --usermode install beamerposter
%           (if the remote repo has moved past your local TL year,
%            point --repository at .../tlnet/historic/systems/texlive/<year>/tlnet-final)
%           Manual route: grab macros/latex/contrib/beamerposter.zip from any
%           CTAN mirror, drop beamerposter.sty into ~/texmf/tex/latex/beamerposter/.
%           Overleaf ships it out of the box.
% Fonts:    Times New Roman (Latin) / SimSun 宋体 (CJK) via fontspec,
%           every optional font guarded by \IfFontExistsTF — missing
%           fonts fall back, never break compilation.
% Canvas:   A0 portrait via beamerposter [size=a0,orientation=portrait];
%           see specs.json for other presets (change size= / orientation=).
% ============================================================
\documentclass[final]{beamer}
\usepackage[size=a0,orientation=portrait,scale=1.0]{beamerposter}

\usepackage{amsmath,amssymb,bm}
\usepackage{fontspec}
\usepackage{xeCJK}              % CJK typesetting (SimSun) for Chinese posters
\usepackage{graphicx}
\usepackage{booktabs}
\usepackage{qrcode}             % \qrcode[height=4cm]{https://doi.org/...}

% ---------- fonts (guarded; Latin Times, CJK SimSun) ----------
\IfFontExistsTF{Times New Roman}{\setmainfont{Times New Roman}}{%
  \IfFontExistsTF{Liberation Serif}{\setmainfont{Liberation Serif}}{%
    \IfFontExistsTF{Nimbus Roman}{\setmainfont{Nimbus Roman}}{}}}
\IfFontExistsTF{SimSun}{\setCJKmainfont{SimSun}}{%
  \IfFontExistsTF{Noto Serif CJK SC}{\setCJKmainfont{Noto Serif CJK SC}}{}}

% ---------- explicit type scale (iron floors from specs.json) ----
% Do NOT tune below: body >= 24pt, captions >= 18pt, heads 36-48pt,
% title 72-85pt at A0. Reading distance: title 3-5 m, body 1-1.5 m.
\newcommand{\fTitle}{\fontsize{76}{88}\selectfont\bfseries}
\newcommand{\fAuthors}{\fontsize{34}{42}\selectfont}
\newcommand{\fHead}{\fontsize{44}{52}\selectfont\bfseries}
\newcommand{\fBody}{\fontsize{26}{34}\selectfont}
\newcommand{\fCaption}{\fontsize{20}{26}\selectfont\itshape}
\newcommand{\fRef}{\fontsize{18}{24}\selectfont}

% ---------- palette (same five colors as slide-making) ----------
\usepackage{xcolor}
\definecolor{cPrimary}{HTML}{0E2841}
\definecolor{cAccent}{HTML}{E97132}
\definecolor{cTeal}{HTML}{156082}
\definecolor{cGray}{HTML}{5A6B7B}
\definecolor{cLight}{HTML}{F2F6F9}

% ---------- beamer look: quiet theme, block styling -------------
\setbeamertemplate{navigation symbols}{}
\usefonttheme{professionalfonts}
\setbeamercolor{structure}{fg=cPrimary}
\setbeamercolor{block title}{fg=white,bg=cPrimary}
\setbeamercolor{block body}{fg=black,bg=cLight}
\setbeamertemplate{blocks}[rounded][shadow=false]

\title{Full Paper Title: Method Subtitle and the Application Domain}
\author{First Author$^{1}$ \and Second Author$^{2}$ \and Third Author$^{1,*}$}
\institute{$^{1}$University Department, City, Country \quad $^{2}$Institute, City, Country\\
           $^{*}$corresponding@university.edu}

\begin{document}

\begin{frame}[t]{}%
\begin{columns}[t]

% ================= LEFT: problem + model =================
\begin{column}{0.31\textwidth}
  {\fTitle \color{cPrimary} Full Paper Title:\\ Method Subtitle\par}
  \vspace{4mm}
  {\fAuthors First Author$^{1}$, Second Author$^{2}$, Third Author$^{1,*}$\par}
  {\fCaption \color{cGray} $^{1}$University \quad $^{2}$Institute \quad
    $^{*}$corresponding@university.edu \quad [CONFERENCE 2026, Poster \#\#]\par}
  \vspace{8mm}

  \begin{block}{\fHead Background \& Problem}
    \fBody
    Two--three sentences of context; one sentence stating the gap:
    existing models fail to handle \emph{[what]}, which matters because
    \emph{[consequence]}. \par\vspace{3mm}
    \begin{center}
      \includegraphics[width=0.92\linewidth]{example-image}% -> figs/problem.pdf (reuse paper Fig. 1)
      {\fCaption \color{cTeal} Figure 1. Problem setting, self-contained caption.\par}
    \end{center}
  \end{block}

  \begin{block}{\fHead Model}
    \fBody Decision problem in words: sets, first-stage decisions,
    recourse, objective. \par\vspace{4mm}
    % Math is why you chose this route — keep the model MINIMAL:
    % objective + two coupling constraints, not the whole appendix.
    \begin{equation*}
      \min_{\bm{x}}\; c^\top \bm{x}
      + \max_{P\in\mathcal{P}} \mathbb{E}_{P}
        \big[ Q(\bm{x},\bm{\xi}) \big]
      \quad \text{s.t.}\;
      A\bm{x} \ge \bm{b},\; \bm{x}\in\mathcal{X}
    \end{equation*}
    \fCaption \color{cGray} Distributionally robust two-stage program;
    ambiguity set $\mathcal{P}=\{P: W(P,\hat P)\le\rho\}$.
  \end{block}
\end{column}

% ================= MIDDLE: method =================
\begin{column}{0.31\textwidth}
  \begin{block}{\fHead Method}
    \fBody Name the framework and the ONE idea that makes it work.
    Algorithm skeleton in 8--12 pseudocode lines, not more. \par\vspace{3mm}
    \begin{center}
      \includegraphics[width=0.95\linewidth]{example-image}% -> figs/framework.pdf (drawio export)
      {\fCaption \color{cTeal} Figure 2. C\&CG solution framework.\par}
    \end{center}
    \fBody Key acceleration: \emph{[e.g.\ Pareto cuts / level sets /
    warm starts]} — this is the novelty claim, say it plainly.
  \end{block}

  \begin{block}{\fHead Experimental Setup}
    \fBody
    \begin{itemize}
      \item Instances: \emph{[family, size]}
      \item Baselines: \emph{[SAA, Benders, heuristic]}
      \item Hardware / solver / limits: \emph{[one line]}
    \end{itemize}
  \end{block}
\end{column}

% ================= RIGHT: results + takeaway =================
\begin{column}{0.31\textwidth}
  \begin{block}{\fHead Results}
    {\fontsize{64}{70}\selectfont\bfseries \color{cAccent} $-18.4\%$\par}
    {\fCaption \color{cGray} expected cost vs.\ best baseline (median, 50 instances)\par}
    \vspace{3mm}
    \begin{center}
      \includegraphics[width=0.95\linewidth]{example-image}% -> figs/pareto.pdf (hero figure, largest on the poster)
      {\fCaption \color{cTeal} Figure 3. Hero result — cost--robustness tradeoff.\par}
    \end{center}
    \fBody Two sentences on what the figure shows.
    \begin{center}
      \begin{tabular}{lr}
        \toprule
        \fBody Method & \fBody Cost \\ \midrule
        \fBody Ours  & \fBody \textbf{100.0} \\
        \fBody SAA   & \fBody 122.6 \\
        \fBody Det.  & \fBody 131.4 \\
        \bottomrule
      \end{tabular}
    \end{center}
  \end{block}

  \begin{alertblock}{\fHead Takeaway}
    \fBody One sentence a visitor can repeat after leaving: what is
    new, and for whom it changes what.
  \end{alertblock}

  \begin{center}
    \qrcode[height=4.5cm]{https://doi.org/10.xxxx/xxxxx}\\[2mm]
    {\fCaption \color{cGray} Scan for paper, data \& code}
  \end{center}

  {\fRef \color{cGray}
    \textbf{References.} Author, A. (2025). Title. \emph{Journal}, 12(3).
    Author, B. (2024). Title. \emph{Proc.\ Conf.}, 1--10.\\
    \textbf{Funding.} Grant \#; no conflict of interest.\par}
\end{column}

\end{columns}
\end{frame}
\end{document}
